\documentclass[12pt]{article}
\usepackage[T1]{fontenc}
\usepackage[utf8]{inputenc}
\usepackage{lmodern}
\usepackage{textcomp}
\usepackage{enumitem}
\usepackage{geometry}
\geometry{margin=1in}
\begin{document}
\begin{center}
Answer Key - Health Sciences - Undergraduate Level Assessment 21 \
\small (Answers are ordered exactly as in the question paper.)
\end{center}

\section*{Multiple Choice}
\begin{enumerate}[label=\arabic*.]
\item \textbf{Answer:} C
\\ \textit{Options:} A) an association.; B) a dysplasia.; C) a sequence.; D) a syndrome.
\\ \textit{Note:} The correct answer is a sequence because pulmonary hypoplasia is a consequence of the primary condition of renal agenesis, which leads to oligohydramnios, demonstrating a chain of developmental abnormalities resulting from a single initial defect.
\item \textbf{Answer:} B
\\ \textit{Options:} A) Puerto Rican; B) Mexican; C) Cuban; D) Brazilian
\\ \textit{Note:} The correct answer is Mexican because the Mexican community constitutes the largest segment of the Latino population in the United States, reflecting historical migration patterns and cultural ties.
\item \textbf{Answer:} D
\\ \textit{Options:} A) Familial adenomatous polyposis; B) Li-Fraumeni syndrome; C) Von Hippel-Lindau syndrome; D) Waardenburg syndrome
\\ \textit{Note:} Waardenburg syndrome is primarily a genetic disorder affecting pigmentation and hearing, rather than a familial cancer syndrome characterized by an increased risk of developing specific types of cancer.
\item \textbf{Answer:} D
\\ \textit{Options:} A) Coping strategies; B) Work productivity; C) Rates of dementia; D) Personality traits
\\ \textit{Note:} Personality traits tend to be influenced more by the unique experiences and cultural contexts of specific cohorts than by the biological aging process itself, making them a key difference observed across different age groups.
\item \textbf{Answer:} B
\\ \textit{Options:} A) Cohort; B) Serial cross-sectional; C) Mortality; D) Syndromic
\\ \textit{Note:} Serial cross-sectional study designs are often employed in surveillance systems because they allow for the collection of data at multiple points in time from different populations, facilitating the monitoring of health trends and prevalence without requiring follow-up of the same individuals.
\end{enumerate}

\section*{True / False}
\begin{enumerate}[label=\arabic*.]
\setcounter{enumi}{5}
\item \textbf{Answer:} False
\\ \textit{Options:} A) True; B) False
\\ \textit{Note:} The statement is false because attraction to similar individuals is influenced by shared attitudes and values, which play a significant role in forming deeper connections beyond mere physical appearance.
\item \textbf{Answer:} True
\\ \textit{Options:} A) True; B) False
\\ \textit{Note:} The needle used in an inferior alveolar nerve block must pass anterior and lateral to the medial pterygoid muscle to effectively target the inferior alveolar nerve while avoiding injury to surrounding structures.
\item \textbf{Answer:} False
\\ \textit{Options:} A) True; B) False
\\ \textit{Note:} The seventh right intercostal space in the mid-axillary line is actually located over the lower lobe of the lung, not the upper lobe.
\item \textbf{Answer:} False
\\ \textit{Options:} A) True; B) False
\\ \textit{Note:} The Old Testament addresses sexuality not only as a biological function but also as a complex social and moral issue, emphasizing the importance of relationships, family structure, and community values.
\item \textbf{Answer:} True
\\ \textit{Options:} A) True; B) False
\\ \textit{Note:} Working in a category IV laboratory, which handles highly pathogenic agents like Ebola, poses the greatest risk for infection due to the high level of biosafety required and the potential for exposure to infectious materials.
\end{enumerate}

\section*{Long Form Response}
\begin{enumerate}[label=\arabic*.]
\setcounter{enumi}{10}
\item \textbf{Answer:} Social Security plays a vital role in ensuring financial stability for retirees, serving as a foundational source of income that significantly influences their overall health and well-being. By providing a guaranteed monthly benefit, Social Security alleviates the financial stress often associated with retirement, allowing individuals to access necessary healthcare services and maintain a decent standard of living. This financial security is crucial for older adults, as it can lead to better health outcomes by reducing anxiety related to economic uncertainty and enabling access to resources that promote wellness. Moreover, Social Security fosters a sense of dignity and independence among retirees, enhancing their quality of life during a critical phase of their lives. Thus, in the context of retirement planning, Social Security is not merely a financial tool but a key component that supports both the economic and physical well-being of retirees.
\\ \textit{Note:} The answer correctly identifies Social Security as a crucial financial support system for retirees, emphasizing its role in reducing financial stress, facilitating access to healthcare, and enhancing overall well-being and quality of life during retirement.
\item \textbf{Answer:} Retirees commonly engage in activities such as traveling, volunteering, and pursuing hobbies like gardening or painting, which promote social interaction and personal fulfillment. These pursuits are often highlighted due to their positive impacts on mental and physical health, as they foster community engagement and lifelong learning. In contrast, home renovation may be less frequently mentioned because it can be more physically demanding and less socially engaging than other activities. Additionally, many retirees may prefer to downsize or simplify their living situations rather than invest time and resources into extensive home projects. Thus, while home renovation can enhance one's living environment, it often does not align with the lifestyle priorities that retirees commonly embrace.
\\ \textit{Note:} correct because it identifies popular retiree activities that enhance social interaction and personal fulfillment, while contrasting them with less commonly mentioned home renovation, which is seen as physically demanding and less socially engaging.
\end{enumerate}

\vfill
\noindent\textit{End of Answer Key}
\end{document}